\chapter{More examples}

\section{Korvai examples}

\paragraph{(A)}
The following Korvai is pretty long. Oddly, it has conclusive fillers everywhere. The
poorva anga is to be recited/played in chatusram, 4 atoms per clap. The entire korvai
is set to Adi tala --- 8 claps per cycle.

\begin{center}
\begin{adjustbox}{max width=\textwidth}
\begin{tabular}{l c}
\toprule
Sol & Claps\\
\midrule
Tha dhee - ki ta thom (6 atoms) & $1\frac{1}{2}$\\
Tham - - - (4 atoms)            & \\
\midrule
Tha ka, tha dhee - ki ta thom (8 atoms) & \multirow{2}{*}{4}\\
Tha ka, tha dhee - ki ta thom (8 atoms) & \\
Tham - - - (4 atoms)                    & \\
\midrule
Tha ka dhi mi, tha dhee - ki ta thom (10 atoms) & \multirow{3}{*}{$7\frac{1}{2}$}\\
Tha ka dhi mi, tha dhee - ki ta thom (10 atoms) & \\
Tha ka dhi mi, tha dhee - ki ta thom (10 atoms) & \\
Tham - - - (4 atoms)                            & \\
\bottomrule
\end{tabular}
\end{adjustbox}
\end{center}

The following is the uttara anga (second part). Recite or play in tisram-double speed,
6 atoms per clap.

\begin{center}
\begin{adjustbox}{max width=\textwidth}
\begin{tabular}{l c}
\toprule
Sol & Claps\\
\midrule
Tha dhee - ki ta thom (6 atoms) & 1\\
Thaa n gu (3 atoms)             & $\frac{1}{2}$\\
\midrule
Tha dhee - ki ta thom (6 atoms)             & \multirow{3}{*}{$2\frac{1}{3}$}\\
(Tha ka) (2 atoms kaarvai - swallowed)      & \\
Tha dhee - ki ta thom (6 atoms)             & \\
Thaa n gu (3 atoms)                         & $\frac{1}{2}$\\
\midrule
Tha dhee - ki ta thom (6 atoms)             & \multirow{5}{*}{$3\frac{2}{3}$}\\
(Tha ka) (2 atoms kaarvai - swallowed)      & \\
Tha dhee - ki ta thom (6 atoms)             & \\
(Tha ka) (2 atoms kaarvai - swallowed)      & \\
Tha dhee - ki ta thom (6 atoms)             & \\
\bottomrule
\end{tabular}
\end{adjustbox}
\end{center}

\paragraph{(B)}
The following korvai is pretty beautiful: the tha dhee - ki ta thom builds up
gradually and concludes in tisram rather than the expected chatusram. The effect is
marvellous. The tala is misra chapu. The walk is chatusram --- 4 atoms per clap ---
except for the last three tha dhee - ki ta thoms. Poorva anga is 10 claps long. Uttara
anga is 11 claps long.

Poorva anga:

\begin{center}
\begin{adjustbox}{max width=\textwidth}
\begin{tabular}{l c}
\toprule
Sol & atoms\\
\midrule
Tha dhi thaka dhina thom - & 6\\
Tha dhi thaka dhina thom - & 6\\
dhi thaka dhina thom -     & 5\\
dhi thaka dhina thom -     & 5\\
thaka dhina thom -         & 4\\
thaka dhina thom -         & 4\\
thaka thom -               & 3\\
thaka thom -               & 3\\
tham -                     & 2\\
tham -                     & 2\\
\bottomrule
\end{tabular}
\end{adjustbox}
\end{center}

Uttara anga:

\begin{center}
\begin{adjustbox}{max width=\textwidth}
\begin{tabular}{l c}
\toprule
Sol & atoms\\
\midrule
Tha - dhee - - - ki - ta - thom - & 12\\
(Tha ka dhi mi)                   & 4\\
Tha dhee - ki ta thom             & 6\\
Tha dhee - ki ta thom             & 6\\
(Tha ka dhi mi)                   & 4\\
\midrule
Tha dhee - ki ta thom (6 atoms) & \multirow{3}{*}{3 claps (Tisram)}\\
Tha dhee - ki ta thom (6 atoms) & \\
Tha dhee - ki ta thom (6 atoms) & \\
\bottomrule
\end{tabular}
\end{adjustbox}
\end{center}

\paragraph{For the still ardent enthusiast}
By retaining everything but the karvai in the uttara anga, one could convert the same
korvai to another tala. Now the karvai is tha ka dhi mi --- 4 atoms long. Changing that
to 8 atoms would add two claps to the entire count. That makes it 24 claps --- Adi and
Rupakam. Changing it to tha ka --- two atoms long, would reduce 1 clap overall and make
it 20 claps long --- Kanda chapu or even Misra-jati Jampa tala --- 7-clap-count +
single-clap + double-clap.

\paragraph{(C)}
The following korvai proves the fact that no matter how small a korvai is, it can still
be superb. Thaka dhin -, tha tha dhin - can be replaced by tha ka dhi mi, tha ki ta or
tha ki ta, tha ka dhi mi.

\begin{center}
\begin{adjustbox}{max width=\textwidth}
\begin{tabular}{l c}
\toprule
Sol & Claps\\
\midrule
Thaka dhin -, tha tha dhin -, tham - (9 atoms)         & \multirow{2}{*}{5}\\
Thaka dhin -, tha tha dhin -, tham -, tham - (11 atoms) & \\
\midrule
Thaka dhin -, tha tha dhin - (7 atoms) & \multirow{3}{*}{$5\frac{1}{4}$}\\
Thaka dhin -, tha tha dhin - (7 atoms) & \\
Thaka dhin -, tha tha dhin - (7 atoms) & \\
\midrule
Tha dhi ki ta thom (5 atoms) & \multirow{5}{*}{$5\frac{3}{4}$}\\
Tham -, tham - (4 atoms)      & \\
Tha dhi ki ta thom (5 atoms) & \\
Tham -, tham - (4 atoms)      & \\
Tha dhi ki ta thom (5 atoms) & \\
\bottomrule
\end{tabular}
\end{adjustbox}
\end{center}

\paragraph{(D)}
This last one is in kanda nadai --- 5 atoms per clap.

\begin{center}
\begin{adjustbox}{max width=\textwidth}
\begin{tabular}{l c}
\toprule
Sol & Claps\\
\midrule
Tha - dhi - tha n gu (7 atoms) & \multirow{2}{*}{$2\frac{2}{5}$}\\
Tha dhi ki ta thom (5 atoms)   & \\
\midrule
Dhi - tha n gu (5 atoms)     & \multirow{2}{*}{2}\\
Tha dhi ki ta thom (5 atoms) & \\
\midrule
Tha n gu (3 atoms) & $\frac{3}{5}$\\
\midrule
Tha dhi ki ta thom (5 atoms) & \multirow{3}{*}{3}\\
Tha dhi ki ta thom (5 atoms) & \\
Tha dhi ki ta thom (5 atoms) & \\
\bottomrule
\end{tabular}
\end{adjustbox}
\end{center}
